\section*{ВВЕДЕНИЕ}
\addcontentsline{toc}{section}{ВВЕДЕНИЕ}

\underline{Актуальность темы}. Современные производственные предприятия с дискретным характером выпуска сталкиваются с необходимостью оперативного составления расписаний, учитывающих загрузку и текущее состояние оборудования. Классические методы целочисленного линейного программирования (MILP) требуют значительных вычислительных ресурсов и времени, что ограничивает их применение в реальном масштабе времени. Развитие методов глубокого обучения открывает возможность построения быстрых нейросетевых диспетчеров, способных адаптироваться к изменяющейся производственной обстановке без полного пересчёта плана. Актуальной является задача создания интеллектуальной системы планирования, объединяющей эвристические алгоритмы и нейросетевые модели для повышения эффективности использования оборудования и сокращения времени опоздания заказов.

\underline{Целью выпускной квалификационной работы} является разработка интеллектуальной системы производственного планирования, обеспечивающей построение выполнимых расписаний и адаптивную диспетчеризацию операций на основе нейросетевых методов.

Для достижения поставленной цели необходимо решить \underline{следующие задачи}:
\begin{enumerate}
    \item Разработать модуль управления заказами и производственными ресурсами с учётом технологических типов операций и связей «ресурс–тип».
    \item Разработать подсистему построения производственных расписаний, использующую эвристический алгоритм EDD+LWKR и нейросетевые диспетчеры (PointerNet, SLIM, PPO).
    \item Реализовать модуль обучения нейросетевых моделей на синтетических данных с сохранением весов и метрик обучения.
    \item Провести сравнительный анализ методов планирования по критериям опоздания, makespan и равномерности загрузки ресурсов.
\end{enumerate}

\underline{Объектом исследования} являются процессы оперативного производственного планирования на предприятиях с гибкой производственной структурой.

\underline{Предметом исследования} являются методы построения расписаний с использованием нейросетевых диспетчеров, эвристических алгоритмов и их комбинаций, а также влияние параметров готовности ресурсов на качество плана.

\underline{Методы исследования}. При проведении исследований использовались методы системного анализа, математического моделирования, эвристической оптимизации, глубокого обучения (нейронные сети прямого распространения, обучение с подкреплением), объектно-ориентированного программирования и сравнительного анализа.

\underline{Научная новизна в практической области}. Разработанная система позволяет в реальном времени формировать производственное расписание, автоматически учитывая текущую загрузку, готовность и специализацию ресурсов. Использование нейросетевых диспетчеров, обученных на эвристически размеченных данных, сокращает время построения плана по сравнению с точными MILP-решателями и обеспечивает возможность адаптации к изменению состояния оборудования (ремонт, снижение надёжности) без переобучения модели.

\underline{Основные результаты}. В ходе исследования и разработки получены следующие результаты:
\begin{enumerate}
    \item Разработан модуль управления заказами и производственными ресурсами, учитывающий технологические типы операций и связи «ресурс–тип».
    \item Разработана подсистема построения производственных расписаний, включающая эвристический алгоритм EDD+LWKR и нейросетевые диспетчеры PointerNet, SLIM, PPO.
    \item Реализован модуль обучения нейросетевых моделей на синтетических данных с сохранением весов и метрик обучения.
    \item Проведён сравнительный анализ методов планирования на стандартизированном тестовом сценарии; установлено, что модель SLIM обеспечивает наилучшие показатели по опозданию, makespan и равномерности загрузки.
\end{enumerate}

\underline{Достоверность полученных результатов} подтверждается применением апробированных методов машинного обучения, воспроизводимостью экспериментов (фиксированное начальное состояние генератора псевдослучайных чисел, единый тестовый сценарий) и согласованностью поведения нейросетевых моделей с эталонной эвристикой.

\underline{Практическая ценность}. Разработка интеллектуальной системы производственного планирования предназначена для автоматизации составления расписаний и оперативной диспетчеризации в цехах с гибкой структурой производства. Система может использоваться как самостоятельный инструмент планирования, так и в составе корпоративной информационной среды предприятия.

\underline{Структура и объем работы}.Отчет состоит из введения, 4 разделов основной части, заключения, списка использованных источников, 2 приложений. Текст отчета выпускной квалификационной работы включает \formbytotal{lastpage}{страниц}{е}{ам}{ам}. % из них 66 страниц основного текста.

\underline{Во введении} обоснована актуальность темы, сформулирована цель работы, поставлены задачи исследования, представлены научная новизна, основные результаты, практическая ценность, описана структура работы.

\underline{В первом разделе} проводится анализ предметной области оперативного производственного планирования. Рассматриваются эволюция систем управления производством, особенности задачи гибкого производственного расписания, существующие ERP, MES и APS-решения и их ограничения. Особое внимание уделяется нейросетевым подходам к диспетчеризации: архитектурам, методам обучения (поведенческое клонирование, самообучение, обучение с подкреплением) и критериям оценки расписаний.

\underline{Во втором разделе} представлены теоретические модели и алгоритмы, лежащие в основе системы. Описывается математическая постановка задачи выбора ресурса, эвристический алгоритм EDD+LWKR с учётом готовности оборудования, архитектура нейросетевого диспетчера на основе многослойного перцептрона и Pointer Network, а также процедуры обучения моделей. Приводятся метрики качества планирования и методы их расчёта.

\underline{В третьем разделе} детально описывается процесс проектирования и разработки программной системы. Излагаются функциональные и нефункциональные требования, архитектура приложения, структура базы данных и пользовательского интерфейса. Приводится спецификация ключевых классов и модулей (APSCore, SimpleDispatcher, Trainer, SelfLabelingTrainer, SimpleAgent), рассматриваются реализованные алгоритмы планирования и обучения, а также процедуры тестирования.

\underline{В четвертом разделе} представлены результаты экспериментального исследования разработанной системы. На стандартизированном тестовом сценарии сравниваются четыре метода планирования: эвристика, PointerNet, SLIM и PPO. Анализируются показатели опоздания, makespan, равномерности загрузки, динамика выполнения заказов, матрицы ошибок классификации ресурсов. Формулируются выводы о преимуществах саморазмеченного обучения и практической применимости системы.

\underline{В заключении} приводятся выводы о достижении цели и выполнении задач практики, излагаются основные результаты, полученные в ходе выполнения работы.

\underline{В приложении А} представлены фрагменты исходного кода программы.

\underline{В приложении Б} представлен графический материал.

